\section{Preliminaries}\label{prelim}
%

In this section we give some basic definitions and results relevant to
the paper, mainly from Cohn's theory of factorization.  Analogous to
integral domains and unique factorization domains in commutative ring
theory, P.M.~Cohn \cite{Cohnfir,Cohnintro} has developed a theory for
noncommutative rings based on the weak algorithm (a noncommutative
generalization of the Euclidean division algorithm) and the notion of
free ideal rings. We present the relevant basic definitions and
results, specialized to the ring $\FX$ of noncommutative polynomials
with coefficients in a (commutative) field $\F$, and also for matrix
rings with entries from $\FX$.

The results about $\FX$ in Cohn's text \cite[Chapter 5]{Cohnfir} are
stated uniformly for algebraically closed fields $\F$. However, those
we discuss hold for any field $\F$ (in particular for $\F_q$ or a
small degree extension of it). The proofs are essentially based on
linear algebra.

%
%
%
%

%
%
%
%
%

Since we will be using Higman's linearization \cite{Hig} to factorize
noncommutative polynomials, we are naturally lead to studying the
factorization of linear matrices in $\fR^{d\times d}$ using Cohn's
theory.

\begin{definition}{\rm\cite{Cohnfir}}
A matrix M in $\fR^{d\times d}$ is called \emph{full} if it has
(noncommutative) rank $d$. That is, it cannot be decomposed as a
matrix product $M= M_1\cdot M_2$, for matrices $M_1 \in \fF^{d \times
  e}$ and $M_2 \in \fF^{e \times d}$ with $e <d$.
\end{definition}

\begin{remark}
 Based on the notion of noncommutative matrix rank \cite{Cohnfir}, the
 square matrix $M\in\fR^{d\times d}$ is full precisely when it is
 invertible in the skew field $\fF$. That is, $M$ is full if and only
 if there is a matrix $N\in\fF^{d\times d}$ such that $MN=NM=I_d$, where $I_d$ is $d \times
d$ identity matrix.
\end{remark}

We note the distinction between full matrices and units in the matrix
ring $\fR^{d\times d}$.

\begin{definition}
A matrix $U \in \fR^{d\times d}$ is a \emph{unit} if there is a matrix
$V\in \fR^{d\times d}$ such that $UV=VU=I_d$, where $I_d$ is $d \times
d$ identity matrix.
\end{definition}

Clearly, units in $\fR^{d \times d}$ are full. Examples of units in
$\fR^{d\times d}$, which have an important role in our factorization
algorithm, are upper (or lower) triangular matrices in $\fR^{d \times
  d}$ whose diagonal entries are all \emph{nonzero scalars}. Full
matrices, in general, need not be units: for example, the $1\times 1$
matrix $x$, where $x$ is a variable, is full but it is not a unit in
the ring $\fR^{1\times 1}=\fR$.

\begin{remark}
Full non-unit matrices are essentially non-unit non-zero-divisors.
For the factorization of elements in $\fR^{d \times d}$, units are
similar to scalars in the factorization of polynomials in polynomial
rings. Cohn's theory \cite{Cohnfir} considers factorizations of full
non-unit elements in $\fR^{d \times d}$.
\end{remark}
%
%
We next define \emph{atoms} in $\fR^{d\times d}$, which are
essentially the irreducible elements in it.

\begin{definition}
A full non-unit element $A$ in $\fR^{d \times d}$ is an \emph{atom} if
$A$ cannot be factorized as $A=A_1A_2$ for full non-unit matrices
$A_1, A_2$ in $\fR^{d \times d}$.
\end{definition}

Noncommutative polynomials do not have unique factorization in the usual
sense of commutative polynomial factorization.\footnote{However, as
  shown by Cohn, using the notion of stable associates there is a more
  general sense in which noncommutative polynomials have ``unique''
  factorization \cite{Cohnfir}.}  A classic example \cite{Cohnfir} is
the polynomial $x+xyx$ with its two different factorizations
\[
x+xyx = x(1+yx) = (1+xy)x,
\]
where $1+xy$ and $1+yx$ are distinct irreducible polynomials.

\begin{definition}
Elements $A \in \fR^{d \times d}$ and $B \in \fR^{d' \times d'}$ are
called \emph{stable associates} if there are positive integers $t$ and
$t'$ such that $d+t=d'+t'$ and units $P, Q\in\fR^{(d+t)\times (d+t)}$
such that $A\oplus I_t=P(B \oplus I_{t'})Q$.
\end{definition}

It is easy to check that the polynomials $1+xy$ and $1+yx$ are stable
associates.

Notice that if $A$ and $B$ are full non-unit matrices that are stable
associates then $A$ is atom if and only if $B$ is atom.  Furthermore,
we note that stable associativity defines an equivalence relation
between full matrices over the ring $\fR$.

We observe that the problem of checking if two polynomials in $\fR$ given as arithmetic formulas 
are stable associates or not has an efficient randomized algorithm
(Lemma~\ref{stable-algo}). 

%
%

%
%
%
%
%
%
%

%
%
%

Now we turn to the problem of noncommutative polynomial factorization.
By Higman's linearization \cite{Hig,Cohnfir}, given a polynomial
$f \in \fR$ there is a positive integer $\ell$ such that $f$ is stably
associated with a \emph{linear matrix} $L \in \fR^{\ell \times \ell}$,
that is to say, the entries of $L$ are affine linear
forms.\footnote{More generally, by Higman's linearization any matrix
  of polynomials $M$ is stably associated with a linear matrix
  $L \in \fR^{\ell \times \ell}$ for some $\ell$.}  Higman's
linearization process is a simple algorithm obtaining the linear
matrix $L$ for a given $f$, and it plays a crucial role in our
factorization algorithm. We first describe it and then recall an
effective version \cite{GGOW20} which gives a simple polynomial-time
algorithm to compute $L$ when $f$ is given as a non-commutative
arithmetic formula. Then we state our stronger result showing that
even if $f$ is given by an algebraic branching program as input
we can compute its Higman linearization in deterministic polynomial
time.

\subsection*{Higman's linearization process}

%
%

We describe a single step of the linearization process. Given an
$m \times m$ matrix $M$ over $\fR$ such that $M[m,m]=f+g\times h$,
apply the following:
\begin{itemize}
\item Expand $M$ to an  $(m+1)\times (m+1)$ matrix by adding a new last
  row and last column with diagonal entry $1$ and remaining new entries
  zero:
\[
\left[
\begin{array}{c|c}
M & 0 \\
\hline
0 & 1
\end{array}
\right]. 
\]
\item Then the bottom right $2\times 2$ submatrix is transformed as
  follows by elementary row and column operations

\[
\left(
\begin{array}{cc}
f+gh & 0 \\
0 & 1
\end{array}
\right)\rightarrow
\left(
\begin{array}{cc}
f+gh & g \\
0 & 1
\end{array}
\right)\rightarrow
\left(
\begin{array}{cc}
f & g \\
-h & 1
\end{array}
\right)
\]

\end{itemize}

Given a polynomial $f\in\FX$ by repeated application of the above step
we will finally obtain a \emph{linear matrix} $L=A_0+\sum_{i=1}^n A_i
x_i$, where each $A_i, 0\le i\le n$ is an $\ell\times \ell$ over $\F$,
for some $\ell$. The following theorem summarizes its properties.

\begin{theorem}[Higman Linearization]{\rm\cite{Cohnfir}}\label{higthm}
  Given a polynomial $f\in\fR$, there are matrices $P, Q\in\fR^{\ell \times \ell}$
  and a linear matrix $L\in\fR^{\ell\times \ell}$ such that
\begin{equation}\label{higeq}
%
\left(
\begin{array}{c|c}
f & 0 \\
\hline
0 & I_{\ell-1}
\end{array}
\right) ~=~PLQ
\end{equation}
with $P$ upper triangular, $Q$ lower triangular, and the diagonal
entries of both $P$ and $Q$ are all $1$'s (hence, $P$ and $Q$ are both
units in $\fR^{\ell \times \ell}$).
\end{theorem}

Instead of a single $f$, we can apply Higman linearization to a matrix
of polynomials $M\in\fR^{m\times m}$ to obtain a linear matrix $L$
that is stably associated to $M$.  We first recall the algorithmic
version of Garg et al.~\cite{GGOW20} in this general form.

\begin{theorem}\label{ehigman}{\rm\cite[Proposition A.2]{GGOW20}}
Let $M \in \fR^{m \times m}$ such that $M_{i,j}$ is computed by a
non-commutative arithmetic formula of size at most $s$ and bit
complexity at most $b$. Then, for $k=O(s)$, in time $\poly(s,b)$ we
can compute the matrices $P, Q$ and $L$ in $\fR^{\ell \times \ell }$
of Higman's linearization such that
\[
\left(
\begin{array}{c|c}
M & 0 \\
\hline
0 & I_{k}
\end{array}
\right) ~=~PLQ.
\], where $\ell= m+k$.
Moreover, the entries of the matrices $P$ and $Q$ as well as $P^{-1}$
and $Q^{-1}$ are given by polynomial-size algebraic branching programs
which can also be obtained in polynomial time.
\end{theorem}

We will sometimes denote the block diagonal matrix
$\left(
\begin{array}{c|c}
M & 0 \\
\hline
0 & I_{k}
\end{array}
\right)$ by $M\oplus I_k$.

We now state our strengthening of Theorem~\ref{ehigman} which enables
us to factorize noncommutative polynomials given as algebraic
branching programs. The complete proof is presented in the appendix.

\begin{theorem}\label{block-higman}
  Let $M \in \fR^{m \times m}$ such that each entry $M_{i,j}$ is a
  polynomial computed by a non-commutative algebraic branching program
  of size at most $s$ and bit complexity at most $b$. Then, for
  $k=O(s)$, in time $\poly(s,b)$ we can compute the matrices $P, Q$
  and $L$ in $\fR^{\ell \times \ell }$ of Higman's linearization such
  that
\[
\left(
\begin{array}{c|c}
M & 0 \\
\hline
0 & I_{k}
\end{array}
\right) ~=~PLQ.
\], where $\ell= m+k$.
Moreover, the entries of the matrices $P$ and $Q$ as well as $P^{-1}$
and $Q^{-1}$ are given by polynomial-size algebraic branching programs
which can also be obtained in polynomial time.
\end{theorem}

%
%
%
%
%

As $P$ and $Q$ are units with diagonal entries all $1$'s, the matrix
$M$ is full iff the linear matrix $L$ is full. Also, the scalar matrix
$M(\overline{0})$ (obtained by setting all variables to zero) is
invertible iff the scalar matrix $L(\overline{0})$, similarly
obtained, is invertible.\\

\subsection*{Invariant Subspaces and Ronyai's Algorithm}

\begin{definition}
Let $A_1, \ldots, A_n \in \mathbb{F}^{d \times d}$. A subspace
$V\subseteq \F^n$ is called as common invariant subspace of $A_1,
\ldots, A_n$ if $A_i v \in V$ for all $i\in [n]$ and $v\in V$. 
\end{definition}

Clearly $0$ and $\F^n$ are, trivially, common invariant subspaces for
any collection of matrices. The algorithmic problem is to find a
\emph{non-trivial} common invariant subspace if one exists. Ronyai
\cite{Ronyai2} gives a randomized polynomial-time algorithm for this problem
when $\F$ is finite field.

\begin{theorem}\label{thm-ronyai}{\rm\cite{Ronyai2}}
Given $A_1, \ldots, A_n \in \mathbb{F}_q^{d \times d}$ there is a
randomized algorithm running in time polynomial in $n,d, \log q$ that
computes with high probability a non-trivial common invariant subspace
of $A_1, \ldots, A_n$ if such a subspace exists, and outputs ``no''
otherwise.
\end{theorem}
  
\begin{remark}
 We should note here, the classical Burnside's theorem \cite{Bur05} for matrix algebras over algebraically closed fields. It essentially shows that the algebra generated by $A_1, A_2, \ldots, A_n$ is the full matrix algebra iff there is no nontrivial common invariant subspace. \end{remark}  
  
\begin{remark}
    As already mentioned in the introduction, Friedl and Ronyai
  \cite{FR85} have shown that over rationals the problem is at least as
  hard as factoring square-free integers, and hence likely to be
  intractable.
  
\end{remark}  


\subsection*{Noncommutative Formulas, Algebraic branching programs} 

Next we recall standard definitions of a noncommutative formulas and
noncommutative algebraic branching programs (ABPs). More details about
noncommutative arithmetic computation can be found in Nisan's work
\cite{N91}:
  
A \emph{noncommutative arithmetic circuit} $C$ over a field $\F$ and
indeterminates $x_1,x_2,\ldots,x_n$ is a directed acyclic graph (DAG)
with each node of indegree zero labeled by a variable or a scalar
constant from $\F$: the indegree $0$ nodes are the input nodes of the
circuit. Internal nodes, representing gates of the circuit, are of
indegree two and are labeled by either a $+$ or a $\times$ (indicating
the gate type). Furthermore, the two inputs to each $\times$ gate are
designated as left and right inputs prescribing the order of gate gate
multiplication. Each internal gate computes a polynomial (by adding or
multiplying its input polynomials), where the polynomial computed at
an input node is just its label.  A special gate of $C$ is the
\emph{output} and the polynomial computed by the circuit $C$ is the
polynomial computed at its output gate. An arithmetic circuit is a
\emph{formula} if the fan-out of every gate is at most one.\smallskip

A noncommutative \emph{algebraic branching program} (ABP) is a layered
directed acyclic graph with one source and one sink. The vertices of
the graph are partitioned into layers numbered from $0$ to $d$, where
edges may only go from layer $i$ to layer $i+1$. The source is the
only vertex at layer $0$ and the sink is the only vertex at layer
$d$. Each edge is labeled with a linear linear form in the
noncommuting variables $x_1, x_2, \ldots, x_n$ The size of the ABP is
the number of vertices. The polynomial in $\F\angle{X}$ computed by
the ABP is defined as follows: the sum over all source-to-sink paths
of the product of the linear forms by which the edges of the path are
labeled.

%
%

\begin{comment}
If instead of homogeneous linear forms we use affine linear forms as
edge labels in the branching program, it will compute a possibly
inhomogeneous polynomial. For more details on noncommutative
arithmetic circuits, formulas, ABPs, the reader is referred to
\cite{Nis91}.
\end{comment}

\section{Some Basic Results}\label{bas-res-sec}

In this section we present some basic results required for our
factorization algorithm.


%

\subsection*{Monic linear matrices}

\begin{definition}\label{def-monic}{\rm\cite{Cohnfir}}
Let $L=A_0+A_1x_1+ \ldots + A_nx_n\in\fR^{d\times d}$ be a linear
matrix, where each $A_i$ is a $d\times d$ scalar matrix over
$\F$. Then $L$ is called \emph{right monic} if the $d\times nd$ scalar
matrix $[A_1~A_2~\ldots~A_n]$ has full row rank. Equivalently, if
there are matrices $B_1, \ldots, B_n \in \mathbb{F}^{d \times d}$ such
that $\Sigma_{i=1}^n A_iB_i = I_d$ (i.e. the matrix
$[A_1~A_2~\ldots~A_n]$ has right inverse).

Similarly, $L$ is \emph{left monic} if the $nd\times d$ matrix $[A_1^T~A_2^T~\ldots~A_n^T]^T$ 
has full column rank. $L$ is called \emph{monic} if it is
both left and right monic.

%
%
%
%
%
%
%
%
\end{definition}

The next two results from Cohn \cite{Cohnfir} are important properties
of monic linear matrices.

\begin{lemma}\label{monic-nonunit}{\rm\cite{Cohnfir}}
A right (or left) monic linear matrix in $\fR^{d\times d}$ is not a unit in
$\fR^{d\times d}$.
\end{lemma}

\begin{proof}
Let $L=A_0 + \sum_{i=1}^n A_i x_i$ be right monic, where each $A_i\in
\F^{d\times d}$. By definition, there are matrices $B_i\in\F^{d\times
  d}, 1\le i\le d$ such that $\sum_{i=1}^d A_iB_i=I_d$. Now, suppose
$L$ is a unit. Then there is a matrix $C\in\fR^{d\times d}$ such that
$CL=I_d$. Let the maximum degree of polynomials occurring in $C$ be
$k$, and let $\hat{C}\in\fR^{d\times d}$ denote the degree $k$
component of $C$ (so each nonzero entry of $\hat{C}$ is a homogeneous
polynomial of degree $k$). Clearly, $\hat{C}\cdot (\sum_{i=1}^n A_i
x_i)=0$. The homogeneity of $\hat{C}$'s entries implies that $\hat{C}
A_i=0$ for each $i$. Hence, $\sum_{i=1}^n \hat{C}A_i B_i =0$ which
implies $\hat{C}=0$, contradicting the assumption that
$C\in\fR^{d\times d}$ is the inverse of $L$. The case when $L$ is left
monic is symmetric.
\end{proof}

Let $f\in\fR$ be a nonzero polynomial and $L$ be a linear matrix
obtained from $f$ by Higman linearization as in Equation~\ref{higeq}.
Clearly, $L$ is a full linear matrix. We show that we can transform
$L$ to obtain a full and right (or left) monic linear matrix $L'$ that is stably
associated to $f$. Furthermore, we can efficiently compute $L'$ and
the related transformation matrices.

%
%

%
%

%
%

%


\begin{theorem}{\rm\cite{Cohnfir}}\label{full-monic}
  Let $L = A_0 + \sum_{i=1}^n A_i x_i$ be a full linear matrix in
  $\fR^{d\times d}$ obtained by Higman linearization from a non
  constant polynomial $f\in\fR$. Then there are deterministic
  $\poly(n,d,\log_2 q)$ time algorithms that compute units $U, U'\in
  \fR^{d\times d}$ and invertible scalar matrices $S,
  S'\in\F_q^{d\times d}$ such that:
\begin{enumerate}  
\item $ULS= L'\oplus I_r$, and $L'$ is right monic. Moreover, if $L$
  is not right monic then $r>0$.
\item $S'LU'=L'\oplus I_{r'}$, and $L'$ is left monic. Moreover, if $L$ is
  not left monic then $r'>0$.
\end{enumerate}  
\end{theorem}

\begin{proof}
We prove only the first part. The second part has an essentially
identical proof.

We present a proof with a polynomial-time algorithm for computing
$L'$. If $L$ is already right monic there is nothing to
show. Otherwise, the row rank of the matrix $B = [A_1~A_2~\cdots~A_n]$
is strictly less than $d$. By row operations we can drive at least one
row of $B$ to zero. So, there is an invertible scalar matrix $U_1\in
\F^{d\times d}$ such that $U_1B$ has its last row as zeros. Now
$U_1A_0$ must have its last row non-zero since $L$ is a full linear
matrix. So the last row of $U_1L$ has only scalar entries and at least
one of these is non-zero. By a column swap applied to $U_1L$ we can
bring this non-zero scalar $\alpha$ in the $(d,d)^{th}$
position. Hence, the $(d,d)^{th}$ entry of $U_1LS_1$ is nonzero, where
$S_1$ is the matrix implementing the column swap. Now, with suitable
row operations using the last row, we can make all entries above the
$(d,d)^{th}$ entry of the $d^{th}$ column zero. Applying column
operations we can make all entries of the $d^{th}$ row to the left of
the $(d,d)^{th}$ entry zero. The resulting matrix is of the form
$RU_1LS_1S' = \tilde{L}\oplus 1$, where the unit $R$ is a linear
matrix and $S'$ is an invertible scalar matrix implementing the row
and column operations.

If $\tilde{L}$ is not right monic, we can recursively apply the above
procedure on $\tilde{L}$ until we finally obtain a unit
$\tilde{U}\in\fR^{d\times d}$ and a scalar invertible matrix
$\tilde{S}\in\F^{d\times d}$ such that $\tilde{U}\tilde{L} \tilde{S}=
L' \oplus I_r$, for some positive integer $r<d$, such that $L'$ is
right monic.

%
%
%
%
%
%
%
%
%
%

%
%
%
%

%
%
%
%

To see why this recursive procedure terminates for $r<d$, note that
the dimension of matrix $\tilde{L}$ is reducing by $1$ in each
recursive step and the matrix $\tilde{L}$ obtained is a stable
associate of $L$. So, if $r=d$ it would imply $L$ is a unit which is a
contradiction as we know that $L$ is obtained via Higman linearization
on a non-constant polynomial $f$, so $L$ is noninvertible.

Putting $U=RU_1\tilde{U}$ and $S=S_1\tilde{S}$ we have $ULS=L'\oplus
I_r$ where $L'$ is right monic as desired. It is clear that the entire
construction is polynomial time bounded, and that we have small ABPs
for the entries of $U$.
\end{proof}

\begin{remark}\label{remark-full-monic}
By repeated application of the algorithm in Theorem \ref{full-monic}
we can compute units $U_1,U_2\in\fR^{d\times d}$ such that $U_1 L U_2
= L' \oplus I_r$, where $L'$ is \emph{both left and right monic}. Such
a two-sided monic $L'$ is called \emph{monic} in \cite{Cohnfir}.

For our factorization algorithm, it suffices to compute an $L'$ that
is either left or right monic that is associated to $L$ as in Theorem
\ref{full-monic}. It turns out that either a left monic or a right
monic $L'$ suffices to use Ronyai's common invariant subspace
algorithm to factorize $L'$ (and hence also $L$) as we show in Theorem
\ref{lfact1}. More importantly, the fact that matrices $S$ and $S'$ in
Theorem~\ref{full-monic} are scalar is important for the factor
extraction algorithm as discussed in Theorem \ref{cnzthm}.
\end{remark}

\begin{lemma}\label{stable-algo}
  Given polynomials $f,g\in\fR$ as input by noncommutative arithmetic
  formulas, we can check in randomized polynomial time if $f$ and $g$
  are stable associates.
\end{lemma}

\begin{proof}
  Given $f$ and $g$, using Higman linearization we first compute in
  polynomial time full and monic linear matrices $A$ and $B$ such that
  $f$ and $A$ are stable associates and $g$ and $B$ are stable
  associates (see Theorem \ref{full-monic} and Remark
  \ref{remark-full-monic}). Now, $f$ and $g$ are stable associates iff
  $A$ and $B$ are stable associates. As both $A$ and $B$ are full and
  monic linear matrices, they are stable associates iff both $A$ and
  $B$ are matrices of the same dimension, say $d$, and there are
  scalar invertible matrices $P$ and $Q$ in $\F^{d\times d}$ such that
  $PA=BQ$ \cite[Theorem 5.8.3]{Cohnfir}, where $\F=\F_q$ or a small field extension. Letting the
  $2d^2$ entries of $P$ and $Q$ be variables, we can find a linearly
  independent set of solutions to $PA=BQ$ in polynomial time. Now,
  there exists invertible $P$ and $Q$ in the solution set iff the
  degree-$2d$ polynomial $\det P \times \det Q$ is nonzero on the
  solutions to $PA=BQ$. We can check this by the
  DeMillo-Lipton-Schwartz-Zippel Lemma \cite{DL78,Sch80,Zip79} by
  evaluating $\det P$ and $\det Q$ on a random linear combination of
  the basis of solutions to $PA=BQ$. This will be correct with high
  probablity.
\end{proof}


The next result shows how irreducibility (more generally, the property
of being an atom) is preserved by Higman linearization.

\begin{theorem}\label{associates-preseve-atoms}
Let $f\in\fR$ be a nonconstant polynomial and $L$ 
%
%
  be a full linear matrix stably associated with $f$ (obtained via
Higman linearization).Then the
polynomial $f$ is irreducible iff $L$ is an atom.
%
\end{theorem}

We give a self-contained proof of the above theorem, using the
following (suitably paraphrased) result of Cohn.

\begin{lemma}[Matrix Product Trivialization]{\rm\cite[pp.~198]{Cohnintro}}\label{triv-lemma}
    Let $A\in\fR^{m\times n}$ and $B\in\fR^{n\times s}$ be polynomial
    matrices such that their product $AB=0$. Then there exists a unit
    $P\in \fR^{n\times n}$ such that for every index $i\in [n]$ either
    the $i^{th}$ column of the matrix product $AP$ is all zeros or the
    $i^{th}$ row of the matrix product $P^{-1}B$ is all zeros.
\end{lemma}    

\begin{proofof}{Theorem~\ref{associates-preseve-atoms}}
By Higman linearization, we have upper and lower triangular matrices
$P$ and $Q$, respectively, such that
\[
f\oplus I_s = PLQ,
\]
for some positive integer $s$.

Now, if $f$ is not irreducible then we can write $f=f_1f_2$, where
$f_1$ and $f_2$ are both nonconstant polynomials in $\fR$. Hence
$f\oplus I_s$ factorizes as the product of non-units $(f_1\oplus
I_s)\cdot (f_2\oplus I_s)$, which implies the factorization
\[
L = P^{-1}(f_1\oplus I_s)(f_2\oplus I_s)Q^{-1}.
\]
Now, we claim $P^{-1}(f_1\oplus I_s)$ and $(f_2\oplus I_s)Q^{-1}$ are
non-units. Suppose $P^{-1}(f_1\oplus I_s)$ is a unit. Then $f_1\oplus
I_s$ is a unit which would imply there is an invertible matrix
$M\in\fR^{d\times d}$ such that $(f_1\oplus I_s)M=I_{s+1}$. But that
implies $f\cdot M_{1,1}=1$ which is impossible since $f_1$ is a
nonconstant polynomial. Similarly, $(f_2\oplus I_s)Q^{-1}$ cannot be a
unit. Hence $L$ is not an atom.

Conversely, suppose $L$ is not an atom. Then we can factorize it as
$L=M_1M_2$, where $M_1,M_2\in \fR^{d\times d}$ are
full non-units. Therefore, we have the factorization
\[
f\oplus I_s = (PM_1)(M_2Q).
\]

Writing the matrices $PM_1$ and $M_2Q$ as $2\times 2$ block matrices,
we have:
\[\left(
\begin{array}{c|c}
f & 0 \\
\hline
0 & I_s
\end{array}
\right) 
= \left(
\begin{array}{c|c}
c_1 & c_3 \\
\hline
c_2 & c_4
\end{array}
\right)\cdot 
\left(
\begin{array}{c|c}
d_1 & d_3 \\
\hline
d_2 & d_4
\end{array}
\right).
\]
{From} the $(2,1)^{th}$ matrix block on the left hand side of the
above equation, we obtain the following matrix identity:
\[
0 = \left(
\begin{array}{c c}
c_2 & c_4
\end{array}
\right)\cdot 
\left(
\begin{array}{c}
d_1 \\
d_2
\end{array}
\right),
\]
where $C=(c_2~c_4)$ is in $\fR^{s\times (s+1)}$ and $D=\left(
\begin{array}{c}
d_1 \\
d_2
\end{array}
\right)$ is in $\fR^{(s+1)\times 1}$. By Lemma~\ref{triv-lemma} there
is a unit $U\in \fR^{(s+1)\times (s+1)}$ such that for every $1\le
i\le s+1$ either the $i^{th}$ column of $C''=C\cdot U$ is all zeros or
the $i^{th}$ row of $D''=U^{-1} D$ is all zeros.  Note that $D''$, and
hence $D$, cannot be the all zeros column as $M_2Q$ is full. So, at
least one entry of $D''$ is nonzero. Hence, at least one column of
$C''$ is all zeros. By a suitable column permutation matrix $\Pi$ we
can ensure that the first column of $C \cdot U \Pi$ is all zeros.
Clearly, first entry of $ \Pi^{-1}U^{-1} D$ is nonzero. Writing $f
\oplus I_s$ as a product of $C'= PM_1U\Pi$ and $D'=\Pi^{-1}U^{-1}M_2Q$
we have
\[\left(
\begin{array}{c|c}
f & 0 \\
\hline
0 & I_s
\end{array}
\right) 
= \left(
\begin{array}{c|c}
c'_1 & c'_3 \\
\hline
c'_2 & c'_4
\end{array}
\right)\cdot
\left(
\begin{array}{c|c}
d'_1 & d'_3 \\
\hline
d'_2 & d'_4
\end{array}
\right),
\]

where $c_2'$ is an all zeros column and $d_1'$ is nonzero. {From} the
$(2,2)^{th}$ matrix block of the above equation, we obtain $c_4' d_4'
= I_s$ so $c_4'$ and $d_4'$ are units. By observing $(2,1)^{th}$
matrix block of the above equation we get $c_4' d_2' =0$, which
implies $d'_2$ is an all zeros column as $c_4'$ is unit.  It follows
that $f=c'_1\cdot d'_1$. Furthermore, it is a nontrivial factorization
because both $c'_1$ and $d'_1$ are non-units (because $C'$ and $D'$
are non-units, and $c'_4$ and $d'_4$ are units).
\end{proofof} 




%
%
%
%
%
%
%
%
%
%
%
%
%
%
%
%
%
%
%
%
%
%
%
%
%
%
%
%
%
%
%
%
%
%

Let $L \in\fR^{d\times d}$ be a full and right (or left) monic linear matrix. Let
$L=A_0+\sum_{i=1}^n A_i x_i$. For a positive integer $\ell$ let $M_i,
i\in [n]$ be $\ell\times \ell$ scalar matrices with entries from $\F$
(or a small degree extension of $\F$). Let $Y_i, i\in [n]$ be
$\ell\times \ell$ matrices whose entries are distinct noncommuting
variables $y_{ijk}, 1\le j, k\le \ell$. Then the evaluation of the
linear matrix $L$ at $x_i\leftarrow Y_i+M_i, 1\le i\le n$ is the
$d\ell\times d\ell$ linear matrix in the $y_{ijk}$ variables:
\[
L' = A_0\otimes I_\ell + \sum_{i=1}^n A_i\otimes M_i +
\sum_{i=1}^n\sum_{j,k=1}^\ell (A_i\otimes E_{jk})\cdot y_{ijk}
\]

\begin{lemma}\label{shift-inv-lemma}
  There is a positive integer $\ell\le 2d$ such that for randomly
  picked $\ell \times \ell$ matrices $M_i, i\in[n]$ ( with entries
  from $\F$ or a small degree extension field) the matrix
  $A_0\otimes I_\ell + \sum_{i=1}^n A_i\otimes M_i$ is an
  invertible matrix.
\end{lemma}

\begin{proof}
Since $L\in\fR^{d\times d}$ is a full linear matrix, it has
noncommutative rank $d$. Hence, by the result of \cite{DM17} for the
generic $2d\times 2d$ matrix substitution $x_i \leftarrow X_i,
i\in[n]$, where $X_i$ is a matrix of distinct commuting variables, the
commutative rank of $L(X_1,X_2,\ldots,X_n)$ is $2d^2$ (which means it
is invertible). Hence there is a least $\ell\le 2d$ such that the
commutative rank of $L(X_1,X_2,\ldots,X_n)$ is $d\ell$, where $X_i$
are generic $\ell\times \ell$ matrices with commuting variables.
Hence, by the DeMillo-Lipton-Schwarz-Zippel lemma \cite{DL78,Sch80,Zip79} the
rank of the scalar matrix $L(M_1,M_2,\ldots,M_n)$ is $d\ell$, where
$M_i$ is a random scalar matrix with entries from $\F$ or a small
extension.
\end{proof}

Finally, we state and prove a \emph{modified version} of a result due
to Cohn that allows us to relate the factorization of a polynomial
$f\in\fR$ to the factorization of its Higman linearization $L$. The
proof is given in the appendix.

\begin{theorem}{\rm\cite[Theorem 5.8.8]{Cohnfir}}\label{cohnthm}
  Let $C\in\fR^{d\times d}$ be a full and right monic (or left monic)
  linear matrix for $d>1$. Then $C$ is not an atom if and only if
  there are $d\times d$ invertible scalar matrices $S$ and $S'$ such
  that
  \begin{equation}\label{eq-cohnthm}
  SCS' =
  \left(
\begin{array}{cc}
A & 0 \\
D & B
\end{array}
\right)
\end{equation}
where $A$ is an $r\times r$ full right (respec. left) monic linear matrix and $B$ is
an $s\times s$ full right (respec. left) monic linear matrix such that $r+s=d$.
\end{theorem}

\begin{remark}
  In \cite{Cohnfir} the theorem is proved under the stronger
  assumption that $C$ is monic. However, as we show, it holds even for
  $C$ that is right monic or left monic with minor changes to Cohn's
  proof. We require the above version for our factorization algorithm.
\end{remark}

%
%
%
%

